\section{Introduction}\label{sec:introduction}

Digital research infrastructures (RIs), cyber-physical systems, and Internet-of-Things (IoT) deployments increasingly require software-defined behavior at the edge, where constrained devices execute application logic under reliability, security, timing, life-cycle, and sustainability constraints. Remote deployments can include sensors, actuators, lightweight gateways, and mixed-criticality nodes that must remain inexpensive and portable while still being supervised continuously. In edge-to-cloud RIs, this supervision is not only an operational requirement but also a prerequisite for reproducible experimentation and future energy or carbon accounting across heterogeneous sites. Traditional embedded deployment practices, including manual flashing, board-specific provisioning, and tightly coupled update channels, are workable at small scale but do not provide fleet-wide orchestration or reliable visibility into which software version is running on each device.

In parallel, cloud-native ecosystems have converged around declarative orchestration, reconciliation loops, immutable artifacts, and service-level observability. Kubernetes has become the de facto orchestration substrate for distributed services because it separates desired state from observed state and continuously drives the latter toward the former through controllers~\cite{kubernetes}. This model has transformed data-center operations, but it also exposes a gap when applied to embedded endpoints. Modern orchestrators assume Linux-capable nodes able to run containers, cluster agents, and rich networking stacks, whereas microcontrollers, sensors, and actuators generally cannot host these components. Nevertheless, cyber-physical deployments still need comparable guarantees for secure attachment, observability, and application life-cycle control.

WebAssembly (WASM) provides a practical execution abstraction for this gap. It supports portable, sandboxed workloads that can run across cloud, fog, and edge environments with a more stable module format than architecture-specific binaries~\cite{haas_wasm}. For embedded systems, WebAssembly Micro Runtime (WAMR) and Zephyr-compatible integrations provide a concrete path toward microcontroller-oriented execution~\cite{wamr,zephyr}. Runtime feasibility alone, however, does not solve orchestration: a cloud-side intent must still be translated into a secure command for a constrained device, and execution evidence must be reflected back into the control plane.

This paper introduces a full-stack framework that addresses this problem through a gateway-centric architecture for the Cloud--Fog--Edge continuum. The framework builds an end-to-end path from Kubernetes CRDs on K3s, through a Rust/Transport Layer Security (TLS) gateway that mediates identity validation and state propagation, to Zephyr firmware running WAMR on embedded-class targets. In this path, the cloud records what should happen, the gateway translates and validates that intent, and the device executes the WASM workload while reporting heartbeat-driven liveness and runtime outcomes. Embedded nodes therefore remain lightweight while still participating in a declarative orchestration model. For research infrastructures, the same path provides stable observation and actuation points that can be reused by sustainability dashboards, energy-aware schedulers, or green experiment-design workflows without moving cloud-native complexity into device firmware.

The work is motivated by a concrete engineering observation: secure enrollment is necessary but insufficient. A practical platform must also define ownership of status updates, avoid false disconnection semantics during transient conditions, and expose enough evidence for operators to trust the control plane. This requirement defines the central novelty of the work: a Kubernetes-native evidence and reconciliation model for non-node embedded WASM endpoints, combining declarative resource modeling, secure gateway mediation, state reconciliation, and deployment reproducibility into a single verifiable operational path.

\subsection{Scope and Contributions}
This work presents a systems-oriented contribution at the intersection of IoT orchestration, embedded computing, and secure cyber-physical infrastructure. It does not claim novelty in WebAssembly, Kubernetes, TLS, or Concise Binary Object Representation (CBOR) as isolated technologies. Instead, it integrates these elements into a coherent platform that closes an implementation gap frequently acknowledged in the literature but less often documented end to end: how declarative cloud-side intent can be made operational on constrained, non-node embedded devices without sacrificing secure attachment, control-plane visibility, and state fidelity.

The contributions are fourfold:
\begin{itemize}[leftmargin=*,nosep]
    \item a Kubernetes-native resource model with CRDs for devices, gateways, and deployable WASM applications (Section~\ref{sec:architecture});
    \item a gateway-mediated TLS/CBOR enrollment and control workflow with explicit identity matching and southbound message translation (Section~\ref{sec:protocol});
    \item a reconciliation strategy that links desired and reported state while avoiding false disconnect loops and ambiguous status ownership (Section~\ref{sec:protocol});
    \item an end-to-end K3s validation study with emulated embedded targets and sustainability-oriented baselines for wire traffic, firmware footprint, and control-plane resource use (Section~\ref{sec:evaluation}).
\end{itemize}
The contribution is positioned as an enabling layer for sustainable digital RIs: it does not yet optimize energy use, but it makes edge-device participation declarative, repeatable, and instrumentable for subsequent power, carbon, and placement studies.

The remainder of the paper is organized as follows. Section~\ref{sec:related-work} reviews related work. Section~\ref{sec:problem} formulates the problem and design requirements; Section~\ref{sec:architecture} presents the architecture and resource model. Section~\ref{sec:protocol} describes the southbound protocol and reconciliation policy; Section~\ref{sec:implementation} summarizes the implementation stack. Section~\ref{sec:evaluation} reports the experimental evaluation. Section~\ref{sec:discussion} discusses implications, validity, and reproducibility. Section~\ref{sec:conclusion} concludes the paper.
